\section{Tests Related to Variances}

Testing variances and dispersions is indispensable not only for evaluating uniformity or risk, but also for validating homoscedasticity assumptions in parametric models.

\begin{teorema}[$\chi^2$ test for the variance of a normal population]
	To test whether the variance of a normal population equals a hypothesized value $H_0: \sigma^2 = \sigma_0^2$, we calculate the sample variance $S^2$ from a sample of size $n$. The test statistic is:
	\begin{align}
		\chi^2 = \frac{(n-1)S^2}{\sigma_0^2}.
	\end{align}
	Under $H_0$, this statistic follows a chi-square distribution with $\nu = n - 1$ degrees of freedom.
	For a two-tailed test ($H_a: \sigma^2 \neq \sigma_0^2$), we reject $H_0$ at level $\alpha$ if $\chi^2 < \chi^2_{1-\alpha/2, n-1}$ or if $\chi^2 > \chi^2_{\alpha/2, n-1}$.
\end{teorema}

\begin{teorema}[$F$-test for comparing two population variances]
	To test the hypothesis of equal variances between two independent normal populations ($H_0: \sigma_1^2 = \sigma_2^2$), we draw samples of sizes $n_1$ and $n_2$ and calculate their sample variances $S_1^2$ and $S_2^2$. The test statistic is the ratio of the sample variances:
	\begin{align}
		F = \frac{S_1^2}{S_2^2}.
	\end{align}
	Under $H_0$, this statistic follows an $F$-distribution of Snedecor with $\nu_1 = n_1 - 1$ numerator degrees of freedom and $\nu_2 = n_2 - 1$ denominator degrees of freedom.
	For a two-tailed test ($H_a: \sigma_1^2 \neq \sigma_2^2$), $H_0$ is rejected if $F > F_{\alpha/2, \nu_1, \nu_2}$ or if $F < F_{1-\alpha/2, \nu_1, \nu_2} = \frac{1}{F_{\alpha/2, \nu_2, \nu_1}}$.
\end{teorema}
